\newpage
\section*{Day 6 exercise --- Bootstrapping}

The bootstrap is a \emph{non-parametric} technique for estimating the variance,
standard error or confidence interval of a statistical estimator when no
analytical expression is available (slides 2--3). The idea is to replace the
unknown underlying distribution by the \textbf{empirical distribution} built
from the observed data and then proceed exactly as we would if the true
distribution were known: simulate many datasets, recompute the estimator, and
read off the variability of the resulting estimates (slides 6--7). Slide
numbers refer to the Day 6 deck \emph{Bootstrapping} (Larsen); the exercises
are Chapter 8 of Ross, \emph{Simulation} (p.\ 152).

\section*{Method}

For a sample $x_1,\dots,x_n$ the empirical distribution function assigns mass
$1/n$ to each observation,
\begin{equation}
    F_e(x) = \frac{1}{n}\sum_{i=1}^{n} \mathbf{1}\{x_i \le x\},
    \label{eq:edf}
\end{equation}
i.e.\ the proportion of observations less than or equal to $x$ (slide 8). Since
$nF_e(x)\sim\mathrm{Binomial}(n,F(x))$, the empirical distribution is unbiased
and consistent,
\begin{equation}
    \mathbb{E}[F_e(x)] = F(x), \qquad
    \mathbb{V}[F_e(x)] = \frac{F(x)\,(1-F(x))}{n},
    \label{eq:edf-moments}
\end{equation}
which justifies using $F_e$ as a surrogate for the unknown $F$ (slide 10).

Drawing an observation from $F_e$ amounts to picking one of the data points at
random: generate a uniform $U$ and return $x_I$ with $I=\mathrm{Int}(nU)+1$
(Ross, p.\ 147), i.e.\ sampling with replacement. The bootstrap algorithm
(slide 11) is then:
\begin{enumerate}
    \item Generate $k$ bootstrap samples, each of $n$ observations drawn with
          replacement from the data.
    \item Compute a bootstrap replicate $\hat\theta^*_i$ for each sample.
    \item Estimate the variance of $\hat\theta$ from the replicates,
\end{enumerate}
\begin{equation}
    \bar\theta^* = \frac{1}{k}\sum_{i=1}^{k}\hat\theta^*_i,
    \qquad
    \hat{\mathbb{V}}[\hat\theta] =
        \frac{1}{k-1}\sum_{i=1}^{k}\bigl(\hat\theta^*_i-\bar\theta^*\bigr)^2 .
    \label{eq:boot-var}
\end{equation}
The two core routines implement exactly this --- resampling with replacement and
the variance of the replicates:
\begin{verbatim}
def bootstrap_replicates(data, statistic, k, rng):
    n = len(data)
    samples = rng.choice(data, size=(k, n), replace=True)
    return np.apply_along_axis(statistic, 1, samples)

def bootstrap_variance(data, statistic, k, rng):
    replicates = bootstrap_replicates(data, statistic, k, rng)
    return statistic(data), replicates.var(ddof=1)
\end{verbatim}

% ================================================================
\section*{Part 1 --- Ross Exercise 13}
% ================================================================

For i.i.d.\ $X_1,\dots,X_n$ with unknown mean $\mu$ and constants $a<b$ we want
\begin{equation}
    p = P\Bigl\{a < \tfrac{1}{n}\textstyle\sum_{i=1}^{n} X_i - \mu < b\Bigr\}.
    \label{eq:p13}
\end{equation}

\textbf{(a) Bootstrap approach.} Under the empirical distribution the mean of
the data is the sample mean $\bar x$, so $\bar x$ plays the role of the
(otherwise unknown) $\mu$. We draw $k$ bootstrap samples, compute each bootstrap
mean $\bar X^*_i$, and estimate $p$ by the proportion that fall within $(a,b)$
of $\bar x$:
\begin{equation}
    \hat p = \frac{1}{k}\sum_{i=1}^{k}
             \mathbf{1}\bigl\{a < \bar X^*_i - \bar x < b\bigr\}.
    \label{eq:phat}
\end{equation}

\textbf{(b) Result.} With $n=10$, data $56,101,78,67,93,87,64,72,80,69$,
$a=-5$, $b=5$ and $k=10{,}000$ bootstrap samples:

\begin{table}[h]
    \centering
    \caption{Part 1 --- bootstrap estimate of $p$ in Equation~\eqref{eq:p13}.}
    \label{tab:part1}
    \renewcommand{\arraystretch}{1.1}
    \begin{tabular}{lr}
        \toprule
        Quantity & Value \\
        \midrule
        $\bar x$  & $76.70$ \\
        $\hat p$  & $0.7658$ \\
        \bottomrule
    \end{tabular}
\end{table}

The estimate $\hat p \approx 0.77$ is the bootstrap probability that the sample
mean of $10$ observations lands within $\pm 5$ of the true mean.
Figure~\ref{fig:part1} shows the bootstrap distribution of $\bar X^* - \bar x$
with the interval $[a,b]$ shaded.

\begin{figure}[h]
    \centering
    \includegraphics[width=0.6\textwidth]{plots/day6/ex13_bootstrap_mean.png}
    \caption{Bootstrap distribution of $\bar X^* - \bar x$ ($k=10{,}000$). The
             shaded band is $[a,b]=[-5,5]$; its mass is $\hat p = 0.7658$.}
    \label{fig:part1}
\end{figure}

% ================================================================
\section*{Part 2 --- Ross Exercise 15}
% ================================================================

We estimate $\mathrm{Var}(S^2)$, the variance of the sample-variance estimator
$S^2=\sum_{i=1}^{n}(X_i-\bar X)^2/(n-1)$, for $n=15$ and data
$5,4,9,6,21,17,11,20,7,10,21,15,13,16,8$. There is no convenient closed form, so
we apply the bootstrap (Equation~\eqref{eq:boot-var} with
$\hat\theta^*_i = S^{2*}_i$): resample with replacement, recompute $S^{2*}$ for
each bootstrap sample, and take the sample variance of the $k=10{,}000$
replicates.

\begin{table}[h]
    \centering
    \caption{Part 2 --- bootstrap estimate of $\mathrm{Var}(S^2)$.}
    \label{tab:part2}
    \renewcommand{\arraystretch}{1.1}
    \begin{tabular}{lr}
        \toprule
        Quantity & Value \\
        \midrule
        Observed $S^2$ & $34.31$ \\
        $\widehat{\mathrm{Var}}(S^2)$ & $57.92$ \\
        \bottomrule
    \end{tabular}
\end{table}

As a validation of the method, the closed-form case of Exercise 14 ($n=2$,
$X=\{1,3\}$, where the four equally likely bootstrap samples give
$S^{2*}\in\{0,2,2,0\}$ and hence $\mathrm{Var}(S^2)=1$) is reproduced by the
program to within simulation error ($\approx 1.00$).
Figure~\ref{fig:part2} shows the bootstrap distribution of $S^{2*}$.

\begin{figure}[h]
    \centering
    \includegraphics[width=0.6\textwidth]{plots/day6/ex15_var_s2.png}
    \caption{Bootstrap distribution of the sample variance $S^{2*}$
             ($n=15$, $k=10{,}000$).}
    \label{fig:part2}
\end{figure}

% ================================================================
\section*{Part 3 --- Sample mean vs.\ sample median}
% ================================================================

The program takes a sample and returns the sample median together with a
bootstrap estimate of its variance, using $k=100$ replicates
(Equation~\eqref{eq:boot-var}). We run it on $n=200$ observations from a
Pareto($\beta=1,k=1.05$) distribution, sampled by inversion $X=\beta U^{-1/k}$
(as on Day 2). We compute the sample mean and median and bootstrap the variance
of each.

\begin{table}[h]
    \centering
    \caption{Part 3 --- bootstrap variance of the mean and the median,
             Pareto($\beta=1,k=1.05$), $n=200$, $k=100$.}
    \label{tab:part3}
    \renewcommand{\arraystretch}{1.1}
    \begin{tabular}{lrrr}
        \toprule
        Estimator & Estimate & Bootstrap variance & Bootstrap std.\ error \\
        \midrule
        Sample mean   & $7.39$ & $8.642$ & $2.940$ \\
        Sample median & $2.06$ & $0.018$ & $0.136$ \\
        \bottomrule
    \end{tabular}
\end{table}

The bootstrap variance of the mean is roughly $470$ times that of the median.
This follows directly from the underlying distribution: for $k=1.05$ the Pareto
mean is finite ($\beta k/(k-1)=21$) but the variance is \emph{infinite}
($k<2$), so the distribution is extremely heavy-tailed --- the largest of the
$200$ observations here is about $678$. A few rare, very large values dominate
the sample mean, making it both imprecise (large bootstrap variance) and biased
low for this sample size ($7.39$ against the true $21$). The sample median, by
contrast, depends only on the central order statistics and is essentially
unaffected by the extreme tail, giving a much smaller and more reliable
variance. This is precisely the ``less simple situation'' (slide 5) for which
the bootstrap is intended, since no simple formula for the variance of the
median is available. Figure~\ref{fig:part3} contrasts the two bootstrap
replicate distributions.

\begin{figure}[h]
    \centering
    \includegraphics[width=0.9\textwidth]{plots/day6/part3_mean_vs_median.png}
    \caption{Bootstrap replicate distributions ($k=100$) of the sample mean
             (left) and sample median (right) for the Pareto sample. The mean
             replicates are far more spread out than the median replicates.}
    \label{fig:part3}
\end{figure}
